% ============================================================================
%  Detekcija klikbejt naslova u sportskim vestima na srpskom jeziku
%  Projekat — Obrada prirodnih jezika 2025/2026 (ETF)
%  Kompajliranje: pdflatex izvestaj && bibtex izvestaj && pdflatex x2
%  (Overleaf: glavni fajl = izvestaj.tex, pdfLaTeX)
% ============================================================================
\documentclass[11pt,a4paper]{article}

\usepackage[T1]{fontenc}
\usepackage[utf8]{inputenc}
\usepackage[serbian]{babel}     % srpski, latinica
\usepackage[margin=2.5cm]{geometry}
\usepackage{graphicx}
\usepackage{booktabs}
\usepackage{array}
\usepackage{caption}
\usepackage{float}
\usepackage{hyperref}
\usepackage{csquotes}

\graphicspath{{figures/}}
\hypersetup{colorlinks=true, linkcolor=black, citecolor=black, urlcolor=blue}

% bezbedno uključi tabelu/sliku i ako fajl još ne postoji
\newcommand{\inputIfExists}[1]{\IfFileExists{#1}{\input{#1}}{\textit{[tabela se generiše: \texttt{#1}]}}}
\newcommand{\figIfExists}[2]{\IfFileExists{figures/#1}{\includegraphics[width=#2]{#1}}{\fbox{\textit{[grafik se generiše: #1]}}}}

\title{\textbf{Detekcija klikbejt naslova u sportskim vestima\\ na srpskom jeziku}}
\author{Filip Nikolić (25/3234) \and Danilo Nikolić (25/3235)}
\date{Obrada prirodnih jezika 2025/2026 --- Elektrotehnički fakultet}

\begin{document}
\maketitle
\begin{abstract}
\noindent
% TODO: 5--8 rečenica — zadatak, domen (sportske vesti, SportKlub), pristup
% (baseline LR/NB, enkoderski BERTić/mBERT, dekoderski ChatGPT), ključni rezultat.
Sažetak rada. Predstavljamo binarnu klasifikaciju klikbejt naspram regularnih
sportskih naslova na srpskom jeziku\dots
\end{abstract}

\tableofcontents
\newpage

% ===========================================================================
\section{Uvod}\label{sec:uvod}
% TODO: zadatak (binarna klasifikacija klikbejt/regularan), motivacija,
% definicija klikbejta u sportskom domenu (kognitivni jaz, preuveličavanje,
% emocionalni naboj, senzacionalističke fraze). Kratak pregled rada.

\subsection{Definicija klikbejta}
% prepisati suštinu iz annotation/guidelines.md

% ===========================================================================
\section{Faza 1 --- Prikupljanje podataka}\label{sec:faza1}
% Izvor iz Faza1-plan.md. Radimo SAMO sa SportKlub (WP REST API).
% TODO: izvor i kriterijumi, tehnika (REST API/scraping), polja metapodataka,
% deduplikacija (egzaktna + near-duplicate), UTF-8, finalni broj naslova.

% ===========================================================================
\section{Faza 2 --- Anotacija podataka}\label{sec:faza2}
% Iz Faza2-plan.md + annotation/guidelines.md.
\subsection{Šema oznaka}
% binarna 1/0 + opcioni pod-tipovi (KJ/PV/EN/SF)
\subsection{Uputstvo za anotaciju}
% definicije + granični slučajevi + primeri (iz guidelines.md)
\subsection{Kalibracija i međuanotatorska saglasnost (IAA)}
% kalibracioni skup (220), nezavisna anotacija, Cohen's kappa iz iaa.py
% TODO: ubaciti finalni kappa i % slaganja
\begin{table}[H]\centering
\caption{Međuanotatorska saglasnost na kalibracionom skupu.}
\begin{tabular}{lr}
\toprule
Mera & Vrednost \\
\midrule
Kalibracionih parova & TODO \\
Procentualno slaganje & TODO\,\% \\
Cohen's kappa & TODO \\
\bottomrule
\end{tabular}
\end{table}
\subsection{Glavna anotacija}
% 50/50 podela, jednostruka anotacija, rešavanje spornih slučajeva

% ===========================================================================
\section{Deskriptivna statistika skupa}\label{sec:stat}
% iz src/annotation/stats.py + src/report/make_figures.py
\subsection{Raspodela klasa i dužina naslova}
\begin{figure}[H]\centering
\figIfExists{raspodela_klasa.png}{0.48\textwidth}\hfill
\figIfExists{duzina_naslova.png}{0.48\textwidth}
\caption{Levo: raspodela klasa. Desno: dužina naslova (reči) po klasi.}
\end{figure}
% TODO: komentar — balans 1100/1100, da li su klikbejt naslovi duži/kraći itd.

% ===========================================================================
\section{Faza 3 --- Obučavanje i evaluacija modela}\label{sec:faza3}

\subsection{Protokol evaluacije}
% 10-slojna stratifikovana CV; nested CV za baseline hiperparametre;
% enkoderski: varijante epoha; dekoderski: ceo skup. Metrike: acc, P/R/F1
% (fokus klikbejt klasa), macro-F1, ROC-AUC.

\subsection{Osnovni modeli (LR + Naivni Bajes)}\label{sec:baseline}
% efekat preprocessinga (lowercasing, stem/lemma) × TF/IDF/TF-IDF × n-gram
\inputIfExists{tables/baseline.tex}
% TODO: analiza — koja preprocessing/feature varijanta najbolje, zašto.

\subsection{Enkoderski LLM (BERTić + mBERT)}\label{sec:encoder}
% mono vs multi; efekat broja epoha (2/3/4)
\inputIfExists{tables/encoder.tex}
% TODO: analiza — mono vs multi, koliko epoha optimalno, poređenje sa baseline.

\subsection{Dekoderski LLM (ChatGPT)}\label{sec:decoder}
% SR vs EN prompt, zero/few-shot, sa/bez definicije; eval na celom skupu
\inputIfExists{tables/decoder.tex}
% TODO: analiza — SR vs EN, zero vs few-shot, diskusija cene/limita API-ja.

\subsection{Uporedni pregled svih modela}\label{sec:poredjenje}
\inputIfExists{tables/poredjenje.tex}
\begin{figure}[H]\centering
\figIfExists{poredjenje_modela.png}{0.7\textwidth}
\caption{F1 (klikbejt klasa) najboljih varijanti po tipu modela.}
\end{figure}
% TODO: glavni nalaz — koji pristup pobeđuje i zašto.

% ===========================================================================
\section{Zaključak}\label{sec:zakljucak}
% šta najbolje radi, ograničenja (domen, jedan izvor, veličina skupa), dalji rad.

% ===========================================================================
\bibliographystyle{plain}
\bibliography{references}

\end{document}
